\chapter{Conclusion}
\label{ch:conclusion}

\section{Summary}
\label{sec:concl-summary}

This thesis asked a controlled question: does converting sentence-level factuality
scores into a summary-level preference through Granular Credit Assignment produce a
more learnable reward-model signal than scoring each candidate summary
holistically. The two conditions were compared under an identical candidate pool,
an identical factuality judge, and an identical Bradley--Terry training recipe, so
that the only manipulated factor was the preference-construction procedure itself.

Averaged over six repeated runs at $n=1{,}000$, GCA preferences yield an
advantage of approximately $+3.08$ percentage points in reward-model pairwise
accuracy. The advantage does not remain equally large as the dataset grows: it is
approximately $-0.42$ percentage points at $n=5{,}000$ and $+0.35$ percentage
points at $n=10{,}000$, while absolute accuracy improves for both conditions.
Uncertainty estimates computed over cross-validation folds are anti-conservative,
because folds within a run are not independent, so the direction and approximate
magnitude of the effect are better supported than any exact significance level.

The experiments therefore do not support a universal claim that GCA is preferable
to holistic factuality supervision. They show that sentence-level aggregation can
improve the learnability of automatically generated pairwise preferences under
particular configurations, most clearly in the $n=1{,}000$ setting. In the
exploratory mode comparison, the advantage appears under evaluator modes that do
not themselves decompose the input, which motivates the hypothesis that external
and evaluator-internal decomposition may act partly as substitutes. At larger
sample sizes the observed advantage is small and inconsistent in sign. Taken
together, the results indicate that the usefulness of granular credit assignment
varies with evaluator configuration, dataset size, and aggregation choice rather
than following from finer granularity as such.

\section{Answers to the Research Questions}
\label{sec:concl-rqs}

\textbf{RQ1} (does GCA produce a more learnable preference signal than holistic
scoring?) is answered \emph{conditionally}. An average advantage of approximately
3 percentage points is repeatedly observed at $n=1{,}000$, favouring GCA in five
of six runs, though the run-level test does not meet the conventional threshold.
It is not observed in the single run at $n=5{,}000$ and is marginal at
$n=10{,}000$. The answer is therefore configuration-dependent rather than
general.

\textbf{RQ2} (how do evaluator configuration and aggregation strategy affect
relative learnability?) is answered \emph{associationally}. Both the aggregation
formula and the evaluator mode are associated with substantial changes in the
outcome, and in the exploratory sweep the two modes in which the evaluator
performs its own internal decomposition are also the two in which GCA shows no
advantage (Chapter~\ref{ch:results}, Section~\ref{sec:results-mode-sweep}).
Because only one mode received repeated-run validation, this pattern is suggestive
rather than established.

\textbf{RQ3} (how stable are the differences across runs and dataset sizes?) is
answered \emph{partially}. The difference is directionally consistent across six
runs at fixed $n=1{,}000$, favouring GCA in five of six. It is not reproduced with
comparable magnitude in the single runs at $n=5{,}000$ and $n=10{,}000$, which use
nested supersets of the same data (Section~\ref{sec:disc-scale}).

\section{Limitations}
\label{sec:concl-limitations}

\begin{enumerate}
  \item \textbf{No human factuality audit.} No independent human annotation was
    collected in the final study, so the evaluator's judgments were never checked
    against human assessments of faithfulness
    (Chapter~\ref{ch:discussion}, Section~\ref{sec:disc-reliability}).
  \item \textbf{The labels are not independent ground truth.} AlignScore supplies
    both the holistic and the sentence-level scores, so it functions as the
    reference standard rather than as a verified one. Every result is conditional
    on that evaluator.
  \item \textbf{Preference learnability is not factual correctness.} The
    dependent variable measures how consistently a preference relation can be
    recovered by the chosen architecture, not whether the preferences are
    factually right (Chapter~\ref{ch:methodology},
    Section~\ref{sec:methodology-scope}).
  \item \textbf{No downstream policy evaluation.} The final study contains no
    DPO or other policy-optimization stage, so it does not establish whether
    models trained on either preference set generate more factually consistent
    summaries.
  \item \textbf{Both reward models are weak in absolute terms.} At $n=1{,}000$
    the conditions reach 0.5295 and 0.5603 pairwise accuracy against a chance
    level of 0.5, so the comparison is between two weak signals
    (Chapter~\ref{ch:discussion}, Section~\ref{sec:disc-weak-signal}).
  \item \textbf{Cross-validation folds are not independent.} Uncertainty
    estimates computed over 30 folds treat correlated observations as
    independent and are therefore anti-conservative; the direction and
    approximate magnitude of the effect are better supported than any exact
    significance level (Section~\ref{sec:disc-stats-caveat}).
  \item \textbf{Incomplete determinism across runs.} Subset selection and
    cross-validation fold assignment are explicitly seeded, but PyTorch's
    generator is not seeded in the cross-validation path, leaving weight
    initialization, batch ordering and dropout uncontrolled
    (Chapter~\ref{ch:implementation}, Section~\ref{sec:impl-rm-training}). Exact
    run-level reproduction from the command-line seed alone is therefore not
    guaranteed.
  \item \textbf{Configuration selected on observed results.} The aggregation
    strategy and evaluator mode were chosen after inspecting $n=1{,}000$ outcomes
    (Chapter~\ref{ch:methodology}, Section~\ref{sec:methodology-exploratory}), and
    the larger-scale subsets are nested supersets of that same data
    (Section~\ref{sec:methodology-nesting}), so no result in this thesis
    evaluates the selected configuration on independent data.
  \item \textbf{Candidate generation confounds temperature with condition.} The
    two candidates differ in decoding temperature, and the lower-temperature
    candidate is preferred substantially more often under both conditions and to
    differing degrees (Chapter~\ref{ch:methodology},
    Section~\ref{sec:methodology-ethics}), so temperature-correlated features may
    carry part of the label information.
  \item \textbf{Source-document truncation.} The reward model receives only a
    truncated prefix of each article, so evidence located later in the source is
    unavailable to it (Chapter~\ref{ch:implementation},
    Section~\ref{sec:impl-rm-training}).
  \item \textbf{Evaluator and dataset specificity.} All findings are obtained on
    CNN/DailyMail with AlignScore, Mistral-7B-Instruct-v0.3 candidates, and a
    RoBERTa-base reward model, and need not transfer to other domains,
    evaluators, or architectures.
  \item \textbf{Sentence segmentation was not evaluated.} Summaries are split
    with a regex-based splitter whose accuracy was never measured
    (Chapter~\ref{ch:implementation}, Section~\ref{sec:impl-preferences}).
  \item \textbf{Nested dataset sizes.} The $n=1{,}000$, $n=5{,}000$ and
    $n=10{,}000$ subsets overlap almost completely
    (Chapter~\ref{ch:methodology}, Section~\ref{sec:methodology-nesting}), so the
    larger-scale runs are not independent samples and no experiment evaluates the
    selected configuration on held-out data.
  \item \textbf{No error-category analysis.} No analysis links the observed
    effects to specific categories of factual error; this was not part of the
    final research questions and is listed under future work.
\end{enumerate}

\section{Future Work}
\label{sec:concl-future}

\begin{itemize}
  \item \textbf{Repeated runs at larger scale.} The $n=5{,}000$ and $n=10{,}000$
    experiments were each executed once. Repeating them would test whether their
    intervals overlap the $n=1{,}000$ estimate, which would indicate run-to-run
    variance, or remain separated, which would support a dataset-size effect
    (Chapter~\ref{ch:discussion}, Section~\ref{sec:disc-scale}).
  \item \textbf{Remove the temperature confound.} Generating both candidates at a
    single decoding temperature, and comparing the resulting preference sets
    against the current ones, would establish how much of the observed structure
    is attributable to temperature-correlated surface features rather than to
    factuality (Section~\ref{sec:disc-confounds}).
  \item \textbf{Measure the effect of source truncation.} Varying the retained
    article length while holding all else fixed would quantify how much of the
    near-chance reward-model accuracy is attributable to evidence falling outside
    the model's input window (Section~\ref{sec:disc-weak-signal}).
  \item \textbf{Test the decomposition-interaction mechanism directly.}
    Section~\ref{sec:disc-interpretation} infers from AlignScore's documented
    behavior that its split modes perform internally much of what GCA performs
    externally. Comparing per-sentence scores under split and unsplit modes on the
    same pairs would test that inference rather than rest on it.
  \item \textbf{Independent validation of the labels.} Human annotation of a
    stratified sample of holistic/GCA disagreement cases, or scoring the same
    candidate pool with a second factuality model, would provide the external
    reference this study lacks and would indicate whether either construction
    procedure produces labels closer to human judgment.
  \item \textbf{Downstream policy evaluation.} Training a summarization policy on
    each preference set and evaluating the factual consistency of its generations
    would test whether differences in preference learnability translate into
    differences in generated output, which this study deliberately does not
    address (Chapter~\ref{ch:methodology}, Section~\ref{sec:methodology-evolution}).
  \item \textbf{Error-category analysis.} Categorizing disagreement cases by
    error type would test whether any advantage concentrates in localized errors
    such as entity or relation mistakes, as originally hypothesized.
\end{itemize}
